\documentclass[10pt,letterpaper]{article}

\usepackage[margin=0.50in]{geometry}
\usepackage[T1]{fontenc}
\usepackage{lmodern}
\usepackage{microtype}
\usepackage{enumitem}
\usepackage{titlesec}
\usepackage[hidelinks]{hyperref}
\usepackage{xcolor}
\usepackage{tabularx}

\input{glyphtounicode}
\pdfgentounicode=1

\pagestyle{empty}
\setlength{\parindent}{0pt}
\setlength{\parskip}{0pt}
\setlength{\tabcolsep}{0pt}
\urlstyle{same}

\definecolor{rulegray}{gray}{0.25}

\titleformat{\section}
  {\large\scshape}
  {}
  {0pt}
  {}
  [\color{rulegray}\titlerule]

\titlespacing*{\section}{0pt}{5.5pt}{2.5pt}

\newcommand{\entry}[4]{%
  \begin{tabularx}{\textwidth}{@{}Xr@{}}
    \textbf{#1} & #2 \\
    \textit{#3} & \textit{#4}
  \end{tabularx}\vspace{-1.5pt}
}

\newcommand{\project}[3]{%
  \begin{tabularx}{\textwidth}{@{}Xr@{}}
    \textbf{#1} \textbar{} #2 & #3
  \end{tabularx}\vspace{-3pt}
}

\newcommand{\resumeitem}[1]{\item #1}

\setlist[itemize]{
  leftmargin=15pt,
  itemsep=0.8pt,
  topsep=1.3pt,
  parsep=0pt,
  partopsep=0pt
}

\begin{document}

\begin{center}
  {\LARGE\textbf{Tyrone Marhguy}}\\[-1pt]
  \textbf{Computer Engineering --- Computer Architecture / RTL / FPGA / ASIC / Verification}\\[1.5pt]
  \href{tel:+12673202734}{+1 (267) 320-2734}
  \;|\;
  \href{mailto:tmarhguy@seas.upenn.edu}{tmarhguy@seas.upenn.edu}
  \;|\;
  \href{https://tmarhguy.com}{tmarhguy.com}
  \;|\;
  \href{https://linkedin.com/in/tmarhguy}{linkedin.com/in/tmarhguy}
  \;|\;
  \href{https://github.com/tmarhguy}{github.com/tmarhguy}
\end{center}

\vspace{-5pt}

\section{Summary}

Computer Engineering junior building from transistor-level CMOS and physical design through RTL, FPGA systems, CPU/GPU microarchitecture, networking, formal verification, and ASIC implementation. Built custom processors, a programmable 3D GPU, low-latency Ethernet/ITCH datapaths, SRAM/MAC blocks, and merged fixes into production EDA toolchains.

\section{Education}

\entry
  {University of Pennsylvania}
  {Philadelphia, PA}
  {B.S.E. Computer Engineering}
  {Aug. 2024 -- May 2028}

\vspace{-2pt}

\begin{itemize}
  \resumeitem{
    Coursework: Computer Organization \& Design, VLSI, Digital Logic, Embedded Systems
  }
\end{itemize}

\section{Experience}

\entry
  {Hardware \& Firmware Engineer}
  {May 2026 -- Aug. 2026}
  {Vero Electric}
  {Philadelphia, PA}

\begin{itemize}
  \resumeitem{
    Developed embedded C/C++ fault-detection firmware and a
    \textbf{BQ79616 16-channel battery-monitoring front end} for a
    \textbf{240 kWh / 120 kW} system; closed PCB ERC with \textbf{0 errors}
    and debugged hardware/firmware bring-up using lab instrumentation.
  }
\end{itemize}

\entry
  {Open-Source EDA Contributor}
  {2026 -- Present}
  {OpenROAD, Verilator, LibreLane, OpenFPGA}
  {C++ / Python / EDA Infrastructure}

\begin{itemize}
  \resumeitem{
    Merged fixes into production EDA tools: repaired an
    \textbf{OpenROAD LEF58/ODB parser}, corrected \textbf{Verilator Linux
    peak-memory accounting}, and restored LibreLane/Yosys synthesis compatibility
    and reporting; LibreLane patches shipped in releases
    \textbf{3.0.8 and 3.0.10}.
  }
\end{itemize}

\entry
  {Software Engineer Intern}
  {May 2026 -- Aug. 2026}
  {Aragorn AI}
  {Remote}

\begin{itemize}
  \resumeitem{
    Shipped production backend services and REST APIs; resolved cross-service
    failures blocking releases and hardened system boundaries through architecture
    and code reviews.
  }
\end{itemize}

\section{Selected Hardware Projects}

\project
  {Tomato --- Custom 32-bit Polymorphic CPU}
  {\href{https://tomato.tmarhguy.com}{Website}
   \textbar{}
   \href{https://github.com/tmarhguy/tomato}{GitHub}}
  {SystemVerilog / FPGA / KiCad / Formal}

\begin{itemize}
  \resumeitem{
    Architected a custom 32-bit CPU around two programmable
    \textbf{LUT3 planes + 8-way carry/flag select}
    (\textbf{524,288 ALU configurations}), a three-source datapath, and
    microcoded control; built a working FPGA computer with assembler and
    TomatoOS over HDMI.
  }

  \resumeitem{
    Verified up to \textbf{130B Verilator vectors} plus formal 1/8/32-bit proofs;
    fabricated the discrete ALU PCB and continued the full 74xx/discrete machine
    toward physical bring-up.
  }
\end{itemize}

\project
  {riscv64xO3 --- 64-bit Superscalar OoO RISC-V}
  {\href{https://github.com/tmarhguy/riscv}{GitHub}}
  {SystemVerilog / RISC-V / AXI4-Lite}

\begin{itemize}
  \resumeitem{
    Built an RV64IM core with A/C bring-up and a
    \textbf{2-wide fetch/decode/issue/commit} pipeline, gshare predictor,
    BTB + RAS, AXI4-Lite memory, and Verilator/cocotb/riscv-tests verification
    with Yosys/OpenLane2 feasibility checks.
  }
\end{itemize}

\project
  {Pineapple GPU P1 --- Standalone Programmable 3D GPU}
  {\href{https://github.com/tmarhguy/PineappleGPU}{GitHub}}
  {Verilog / Artix-7 / DVI}

\begin{itemize}
  \resumeitem{
    Implemented programmable vertex/fragment shading, triangle rasterization,
    depth testing, texture sampling, and double-buffered \textbf{320x180} output;
    rendered interactive pineapple and cube scenes on a Nexys A7 using a
    FOSS Yosys/nextpnr flow.
  }
\end{itemize}

\project
  {Low-Latency FPGA Networking \& NASDAQ ITCH}
  {\href{https://github.com/tmarhguy/udp-stack}{UDP/IP}
   \textbar{}
   \href{https://github.com/tmarhguy/itch-hw}{ITCH}}
  {SystemVerilog / RMII / cocotb}

\begin{itemize}
  \resumeitem{
    Built a deterministic \textbf{100 Mbps RMII/MAC/IPv4/UDP} pipeline and an
    RTL ITCH 5.0 parser with a 512-entry BRAM limit book; closed networking at
    \textbf{100 MHz} and verified ITCH against a Python/cocotb golden model over
    a \textbf{10M-message} regression with zero BBO mismatches.
  }
\end{itemize}

\project
  {ASIC / Full-Custom Datapaths}
  {\href{https://github.com/tmarhguy/mac}{MAC}
   \textbar{}
   \href{https://github.com/tmarhguy/64b-sram}{SRAM}}
  {Sky130 / 22 nm HP / SPICE}

\begin{itemize}
  \resumeitem{
    Designed a BF16 multiply-to-FP32-accumulate MAC with LibreLane/OpenLane
    hardening and a transistor-level \textbf{16x4 6T SRAM} with StrongARM sensing;
    automated NGSpice sweeps and achieved simulated functional readback to
    \textbf{4.571 GHz}.
  }
\end{itemize}

\section{Technical Skills}

\textbf{Architecture / RTL:}
SystemVerilog/Verilog, VHDL, CPU/GPU/SoC microarchitecture, pipelining,
superscalar/OoO design, ISA/datapath design, branch prediction,
memory systems, CDC \\

\textbf{Verification / ASIC:}
UVM, SVA, cocotb, formal, Verilator; Vivado, Yosys, OpenROAD,
LibreLane/OpenLane, Sky130, synthesis, STA, P\&R, DRC/LVS, SPICE,
KiCad, Electric VLSI \\

\textbf{Programming / Interfaces:}
Python, C/C++, Embedded C, Tcl, Bash, Linux;
Ethernet/RMII/IPv4/UDP, AXI4-Lite, UART, SPI, I2C

\end{document}